\documentclass[11pt]{article}
\usepackage{preamble}
\usepackage{multicol}
\setlength{\parskip}{4pt}

\begin{document}

\daybanner{AP Statistics \textperiodcentered\ Unit 1 \textperiodcentered\ Topics 1.1--1.2 (CED 1.1--1.2) \textperiodcentered\ TEACHER KEY}{Tile Thickness: Which Records Belong?}

\sectionbanner{CED TETHER}
\begin{multicols}{2}
{\footnotesize\setlength{\parskip}{0pt}
\begin{itemize}[leftmargin=*,itemsep=0pt,topsep=0pt]
\item 1.A: Identify individuals, variables, units, and variation in a statistical context.
\item 2.A: Distinguish categorical labels from quantitative measurements.
\end{itemize}}
\end{multicols}
\clearpage
\small
\textbf{Essential question:} When does a changed percentage answer a different question?\par
\textbf{DOK-3 skill:} 4.B \textperiodcentered\ \textbf{FRQ pattern:} {\small\texttt{defend-\allowbreak{}a-\allowbreak{}tile-\allowbreak{}inspection-\allowbreak{}report}}\par
\textbf{DOK rationale:} Part (d) requires choosing and defending a reporting policy against a stated question, reconciling two correct percentages, categorical codes, and observed variation. Arithmetic alone cannot justify which records belong.\par

\textbf{Self-paced bonus sheet.} Students take it when they choose and turn it in when they finish. Everything they need is printed on the sheet. Score part (d) only, E / P / I, by hand --- never AI-graded or auto-scored. (a) and (b) and (c) are the ladder up; use them to see where a thin (d) came from.\par

\textbf{Questions to ask}
\begin{itemize}[leftmargin=1.5em,itemsep=2pt,topsep=2pt]
  \item What does each numeric column represent?
  \item Which records belong to the original question?
  \item What changes besides the qualifying count when records are omitted?
\end{itemize}

\begin{callout}{\IconWarn\ LOOK FOR}{calloutred}
\begin{itemize}[leftmargin=1.5em,itemsep=2pt,topsep=2pt]
  \item Reporting a percentage without identifying tiles and the thickness threshold.
  \item Treating stage codes or ZIP codes as measured amounts.
  \item Changing only the numerator after removing records.
  \item Confusing differing measured thicknesses with an established cause of variation.
\end{itemize}
\end{callout}

\sectionbanner{SCORING PART (d) --- E / P / I}

\textbf{Expected elements}
\begin{itemize}[leftmargin=1.5em,itemsep=2pt,topsep=2pt]
  \item \textbf{e1} (required) --- Interprets $33.3\%$ and $20.0\%$ as shares of the appropriate tile sets above 8 mm, and compares the full-set result with $25\%$.
  \item \textbf{e2} (required) --- Defends a report answering the original question using the all-valid-record inclusion rule; explains why B changes the set described.
  \item \textbf{e3} (required) --- Explains thickness variation as different measured amounts across tiles and rejects variation alone as grounds for deleting valid records.
  \item \textbf{e4} (required) --- Explains that stage codes label production stages and ZIP codes label postal areas, so averaging either does not answer the thickness question.
\end{itemize}

{\small\renewcommand{\arraystretch}{1.25}
\noindent\begin{tabular}{|p{0.5in}|p{5.9in}|}
\hline
\textbf{E} & All four elements with a coherent defended recommendation. Teacher scores part (d) by hand. \\ \hline
\textbf{P} & A defended recommendation with two or three elements, or all four with a minor calculation or contextual omission. \\ \hline
\textbf{I} & Zero or one element, or an unsupported recommendation with no defensible connection to the original question. \\ \hline
\end{tabular}}

\clearpage
\normalsize
\focusbanner{THE PROBLEM --- annotated}

Suppose a ceramics cooperative inspects 12 tiles. The question is: \emph{Do more than $25\%$ of these 12 inspected tiles have thickness greater than 8 millimeters (mm)?} All 12 measurements are valid, and the inspection rule includes every inspected tile.\par\smallskip
Stage codes mean $1=$ formed, $2=$ bisque-fired, and $3=$ glazed. A studio ZIP code identifies a postal area. Each table row below describes one tile; repeated values represent distinct tiles.
\begin{center}\begin{tabular}{ccc}\hline
\textbf{Stage code} & \textbf{Studio ZIP code} & \textbf{Thickness (mm)} \\\hline
1 & 02139 & 6 \\
1 & 02139 & 7 \\
1 & 02139 & 7 \\
1 & 02139 & 9 \\
2 & 05301 & 6 \\
2 & 05301 & 7 \\
2 & 05301 & 8 \\
2 & 05301 & 9 \\
3 & 02139 & 7 \\
3 & 05301 & 8 \\
3 & 02139 & 10 \\
3 & 05301 & 11 \\\hline
\end{tabular}\end{center}
\textbf{Proposal A:} Keep all 12 records and report the share above 8 mm.\par
\textbf{Proposal B:} Omit the 10-mm and 11-mm tiles as unusually thick, report the remaining share, and announce that the cooperative meets the stated inspection criterion. No measurement error was found.\par
\begin{firsttakebox}{FIRST TAKE --- one sentence before you work the parts}
Which proposal would you provisionally use to answer the inspection question? Give one reason.\par
\answer{Ungraded initial judgment; accept a position with a reason.}
\end{firsttakebox}

\begin{callout}{\IconBook\ NOTES: CONTEXT, LABELS, AND PERCENTS}{calloutgreen}
Individuals are the objects described by the data; variables are recorded characteristics. Quantitative variables measure amounts with units. Categorical variables identify groups; numeric codes may label ordered stages or places without measuring amounts. Variation means that a variable takes different values across individuals.\par A contextual report names the individuals, variable, units, and question. A percentage is $100\times\text{qualifying count}/\text{total count}$; percentage-point change is new percent minus old percent. State the denominator and disclose omitted records. A reporting decision should address the stated question and inclusion rule.
\end{callout}

\dokbadge{1}~\textbf{(a)} Identify the individuals, the measured thickness variable, and its unit. State in words what counts as exceeding the threshold. Checklist: object; characteristic; unit; comparison boundary.\par
\answer{Individuals are the 12 inspected tiles. Thickness is measured in millimeters. A tile exceeds the threshold when its thickness is strictly greater than 8 mm; a tile at 8 mm does not qualify.}

\dokbadge{2}~\textbf{(b)} Classify stage code, studio ZIP code, and thickness as categorical or quantitative. Explain what each code represents and why numerical appearance alone does not justify averaging it. Explain what variation in thickness means here, citing two table values. Checklist: three classifications; meanings of both codes; two measurements; contextual explanation.\par
\answer{Stage and ZIP codes are categorical. Stage labels formed, bisque-fired, and glazed are ordered production stages, not measured amounts. ZIP codes label postal areas, not distances or quantities. Thickness is quantitative in millimeters. Thickness varies because different tiles have different measured thicknesses, such as 6 mm and 11 mm. These observations alone do not establish why the physical differences arose.}

\dokbadge{2}~\textbf{(c)} Calculate the percentage above 8 mm under each proposal, rounded to one decimal place. Give the change from A to B in percentage points and interpret both percentages in context. Checklist: qualifying counts; denominators; calculations; units; which tiles each percent describes.\par
\answer{A: $4/12\times100=33.3\%$. B: $2/10\times100=20.0\%$. B minus A is about $-13.3$ percentage points. The first is the percentage of all inspected tiles thicker than 8 mm; the second is the percentage among the ten retained tiles. Both the qualifying count and denominator change.}

\dokbadge{3}~\textbf{(d)} $\star$ Recommend A, B, or a revised report to answer the original inspection question. Defend your decision with both percentages and the inclusion rule. Address the objection that unusually thick tiles should be removed merely because their thicknesses vary. Explain why averaging stage or ZIP codes would not resolve the decision. Checklist: recommendation; threshold comparison; record inclusion; variation; meanings of both codes. The optional frame below supports only the code explanation.\par
\answer{Use A for the original question: $33.3\%$ of the 12 inspected tiles exceed 8 mm, above the $25\%$ boundary. B correctly calculates $20.0\%$ for ten retained tiles but answers a different question. The two omitted tiles have valid measurements and belong under the inclusion rule, so ordinary observed variation does not justify removing them. Stage codes label production stages and ZIP codes label postal areas; neither numerical average is an amount relevant to the thickness criterion. A revised report may show both results if it identifies the omissions and retains the all-tile conclusion.}

\begin{wordbankbox}\raggedright
\textbf{Word bank:} \textbf{quantitative} \ensuremath{\;\cdot\;} \textbf{postal areas} \ensuremath{\;\cdot\;} distances \ensuremath{\;\cdot\;} \textbf{production stages} \ensuremath{\;\cdot\;} \textbf{millimeters} \ensuremath{\;\cdot\;} percentages \ensuremath{\;\cdot\;} kilograms \ensuremath{\;\cdot\;} \textbf{categorical}
\par\smallskip{\footnotesize\color{framegray} bold = needed; the rest are distractors}
\end{wordbankbox}

\end{document}
